\label{chap:objetivos}
Los objetivos se mantienen tal como fueron aprobados en el protocolo del trabajo de graduación. Las decisiones técnicas que durante el desarrollo modificaron la forma de alcanzarlos, como el cambio de librería de escaneo o de almacenamiento, se explican y justifican en el capítulo de desarrollo del prototipo.

\section{Objetivo general}

Reducir el riesgo de acceso a sitios maliciosos codificados en códigos QR por usuarios móviles, mediante el desarrollo de una interfaz multiplataforma en React Native que integra el flujo completo de escaneo, presentación de resultados y consulta a las APIs de threat intelligence del primer filtro.

\section{Objetivos específicos}

\begin{enumerate}
	\item Implementar el módulo de escaneo de códigos QR con extracción de URL en menos de 5 segundos y clasificación del tipo de contenido, mediante \texttt{react-native-vision-camera} con CodeScanner nativo en iOS y Android.
	\item Desarrollar el flujo de interacción principal entre el escaneo y la presentación del veredicto, mediante la implementación de la pantalla de carga con animación skeleton/shimmer y la pantalla de resultado con dos niveles de detalle (simple y extendido).
	\item Implementar las funcionalidades secundarias de historial offline, reporte y alerta social, mediante almacenamiento local sin cuenta de usuario, integración con el endpoint de reportes y generación de Story Card con geocodificación inversa.
	\item Integrar el filtro F1 de consulta a listas de reputación con tiempo de respuesta inferior a 5 segundos, mediante la conexión con las APIs de Google Safe Browsing, PhishTank y VirusTotal como primer filtro de análisis.
\end{enumerate}
